\def\level{BTS}  % Classe à droite
\def\course{CIEL 2}
\def\chaptername{Lois continues} % Titre au centre
\def\docpart{Loi uniforme~--~TP1} % Partie du cours
\def\docname{C226}        % Identifiant à gauche

\pagestyle{cours_FP}
\makeheaderBTS{} % Affiche le bandeau

\begin{exercice_cours}[\quad]{}

$U$ est une variable aléatoire qui suit la loi uniforme sur l'intervalle $\ff{-2}{8}$.
	\begin{enumerate}
		\item Déterminer la densité de probabilité de $U$.
		\item Calculer $P(U\leqslant 3)$, 
		\item Calculer $P(2\leqslant U\leqslant 5)$ 
		\item Calculer la probabilité d'obtenir une valeur de $U$ inférieure à 5.
		\item Calculer la probabilité d'obtenir une valeur de $U$ inférieure à 3 sachant qu'elle est supérieure à 2.
	\end{enumerate}
\end{exercice_cours}

\begin{exercice_cours}[\quad]{}

	Le temps d'attente à un arrêt de bus peut être modélisé par une variable aléatoire $T$ qui suit le loi uniforme sur un intervalle $\ff{1}{9}$.

	\begin{enumerate}
		\item Déterminer la densité de probabilité de $T$.
		\item Calculer $P(T\leqslant 3)$, 
		\item Calculer $P(2\leqslant T\leqslant 5)$ 
		\item Calculer la probabilité d'attendre moins de 5 minutes.
		\item Calculer la probabilité d'attendre exactement 3 minutes.
		\item Calculer l'espérance de $T$ et donner une interprétation de ce résultat.
	\end{enumerate}


\end{exercice_cours}


\begin{exercice_cours}[\quad]{}

On suppose que la variable aléatoire $X$ suit la loi uniforme sur l'intervalle $\ff{2}{b}$, avec $b>2$.

\begin{enumerate}
	\item Déterminer la densité de probabilité de $X$.
	\item Déterminer $b$ sachant que $P(X\leqslant 3)=\frac{1}{2}$.
	\item Déterminer $b$ sachant que $P(X>8)=\frac{1}{4}$.
\end{enumerate}
\end{exercice_cours}

\begin{exercice_cours}[\quad]{}

	Lors de la restitution des notes d’une interrogation notée sur
20, le professeur annonce que suite à de nombreuses tricheries,
il a décidé de noter chaque élève aléatoirement entre 5
et 15.

\begin{enumerate}
	\item Quelle est la probabilité qu’un élève ait un note supérieur
à 12 ?
	\item Quelle note (moyenne) un élève peut-il espérer obtenir
?
	\item Un élève qui n’a pas triché est effondré à l’annonce de
ce système de notation. Il espérait avoir une note supérieure
à 14. Pour le rassurer, le professeur lui indique
que sa note est supérieur à 12. Quelle est la probabilité
que sa note soit supérieur à 14 sachant qu’elle est supérieur
à 12.
\end{enumerate}
	\end{exercice_cours}

